% !PRINTED 439
\weberpage{439}
%<<CONT
also
\[ 4 p^\varepsilon = y^2 + m x^2. \]

Daraus sind $x$ und $y$ bis auf die Vorzeichen bestimmt, und $x$ kann positiv angenommen werden, und das Vorzeichen von $y$ ergibt sich aus der Kongruenz
\weq{7}{y \equiv b x \pmod{p^\varepsilon}.}

Betrachtet man dagegen die Komposition $l^{-1} k$, so hat man die Form $(a, b, c p^\varepsilon)$ mit $(p, -b, a c p^{\varepsilon - 1})$ zu komponieren; man bestimmt $b'$ aus den beiden Kongruenzen:
\[ b' \equiv b \pmod{2a}, \qquad b' \equiv -b \pmod{2 p^\varepsilon}, \]
und es ist alles ebenso durchzuführen, nur daß an Stelle von (7) der Kongruenz % [sic] printed "der Kongruenz"; grammatically expected "die Kongruenz"
\weq{8}{y \equiv -b x \pmod{p^\varepsilon}}
tritt, d. h. $y$ bekommt das entgegengesetzte Zeichen.

Nach (5) und (6) ist aber
\weq{9}{\alpha + \beta\omega = \frac{y + (2 a \omega + b) x}{2} = \frac{y + \sqrt{-m}\, x}{2},}
worin, wenn $a$ und $x$ und folglich $\beta$ positiv angenommen sind, $\sqrt{-m}$ positiv imaginär zu nehmen ist.

Es ist aber nach \S\,38, (4)
\[ \left( \frac{\eta(\omega_\varepsilon)}{\eta(\omega)} \right)^{24} = \left( \frac{y + \sqrt{-m}\, x}{2} \right)^{12}, \]
und folglich erhält man nach (3)
\weq{10}{\Phi(k, lk) = \left( \frac{y + \sqrt{-m}\, x}{2} \right)^{12}.}

Auf die gleiche Weise ergibt sich wegen (16): % [sic] printed (16); no antecedent equation (16) occurs in this section here
\weq{11}{\Phi(k, l^{-1}k) = \left( \frac{y - \sqrt{-m}\, x}{2} \right)^{12},}
und die Gleichung (10) ist daher \gesperrt{nicht erfüllt}, wenn $(lk)$ durch $(l^{-1}k)$ ersetzt wird, außer wenn
\weq{12}{\left( \frac{y + \sqrt{-m}\, x}{2} \right)^{12} = \left( \frac{y - \sqrt{-m}\, x}{2} \right)^{12}}
ist.

Dies ist aber nicht möglich, wenn, wie wir angenommen haben, $p$ nicht in $m$ aufgeht. Denn zerlegt man $p$ im Körper $\Omega = \Re(\sqrt{D})$ in die zwei voneinander verschiedenen konjugierten Primideale $\mathfrak{p}, \mathfrak{p}'$, so geht das eine in $\frac{1}{2}(y + \sqrt{-m}\, x)$, das andere
%CONT>>
